\section{Complete positive-residence persistence without comparison corridors}
\label{sec:v105-persistence}
Throughout, $(a(\xi),b(\xi))$ is the actual maximal regular roof
segment on a disjoint physical pairing sheet,
$\ell(\xi)=b(\xi)-a(\xi)$, and $c^\varepsilon(\xi)$ is its
unchanged positive endpoint-trace minimum. All exact labels,
the original section and the original event are retained.
For an arbitrary Borel selector $0\le H(\xi,t)\le1$ set
\[
 T_{H,[h,\Delta)}^\varepsilon(t)=
 \int c^\varepsilon(\xi)H(\xi,t)
 \mathbf1_{\{h\le\ell(\xi)<\Delta\}}
 \mathbf1_{\{a(\xi)<t<b(\xi)\}}\,d\theta(\xi).
\]
\begin{theorem}[Original source persistence at every positive residence]
\label{thm:v105-persistence}
At each finite collision count, for every $0<h<\Delta$ and almost
every original roof $t$,
\begin{equation}\label{eq:v105-persistence}
0\le T_{H,[h,\Delta)}^\varepsilon(t)
\le {1\over h}\int_{t-h}^{t+h}\mathsf T_\Delta^\varepsilon(u)\,du
=2A_h\mathsf T_\Delta^\varepsilon(t).
\end{equation}
Consequently, for every previously fixed positive $h,\Delta,\varepsilon$,
\begin{equation}\label{eq:v105-height}
\mathcal H_M(T_{H,[h,\Delta)}^\varepsilon)
\le C\{\varepsilon^{1/16}+(1+\varepsilon^{-1})\Delta\}.
\end{equation}
The leading constant is independent of the previously fixed $h$.
The conclusion covers all physical endpoint types within the
inherited paired source: regular and nonstationary pairs, flat
endpoint pulses, reach-deficient arms and singular extensions.
No positive interval comparison or endpoint derivative is required.
\end{theorem}
\begin{proof}
If $a<t<b$ and $\ell=b-a\ge h$, put
$d_-=t-a$ and $d_+=b-t$. Then
\[
\int_{t-h}^{t+h}\mathbf1_{\{a<u<b\}}\,du
=\min(h,d_-)+\min(h,d_+)\ge h
\]
because $d_-+d_+\ge h$. Multiply this scalar inequality by
$c^\varepsilon(\xi)H(\xi,t)$ and sum over the actual Borel
physical sheets. The selector remains at the original target roof
and is neither differentiated nor transported. Tonelli permits
enlargement of the positive right-hand source to the complete
$\mathsf T_\Delta^\varepsilon$. This proves the finite-count inequality
before taking any collision limit. Equation
\eqref{eq:v104-paired-average}, used only with the fixed
window $h$, gives the ordered essential-height estimate.
No multiplicity or angular word count enters the argument.
\end{proof}
\begin{corollary}[Physical flat pulses at positive residence]
\label{cor:v105-flat}
For every $B$ choose $h_B\in(0,\delta(B))$ from the original
outer parameters, before the collision limit. The selected
original source on all regular pairing segments
$h_B\le\ell\le\delta(B)$, including every positive-residence
flat-pulse sheet even when no comparison interval exists, has
ordered height $O(B^{-1/192})$, uniformly in the choice of
previously fixed $h_B$.
\end{corollary}
\begin{proof}
Use $\Delta=4\delta(B)$ in Theorem
\ref{thm:v105-persistence}, and the original schedule
$\delta(B)=O(\varepsilon(B)^4)$,
$\varepsilon(B)=A_0B^{-1/12}$. The first exponent is
$B^{-1/192}$ and the second contribution decays faster.
\end{proof}
The quantifier is important. A leading constant uniform over
positive \emph{fixed} windows does not imply a finite-count
estimate at $h=0$, nor may a window shrink with collision count.
The complete microscopic residence source remains positive
and is isolated below.
